\chapter{Native fields and the boundary action}
\label{ch:vol2-native-bulk-boundary}

A boundary current can be obtained by integrating fields in a normal
direction. The integration also determines how an interior insertion acts
on that current. These two outputs contain different information.
An eliminated field may have zero projection and still contract nontrivially
with a boundary observable. The resulting operation depends on the
propagator, its supports and its quantum normalization.

We construct this comparison for the Gaussian theory on a three-dimensional
half-space. Its Wilson source then gives an independent incoming junction,
including a charged vacuum and a sewn operation. Finite regular jets can be
retained by explicit changes of representatives. The completions used for
these constructions will be specified separately: a smooth smearing, a
polynomial state and a formal series of states have different pole bounds.

\section{The half-space action}

Put \(M=\mathbb R_{\geq0,t}\times\mathbb C_z\), with
\(z=x+iy\) and orientation \(dt\wedge dx\wedge dy\).
The partial forms on \(M\) are generated by \(dt,d\bar z\). Write
\[
 D=dt\,\partial_t+d\bar z\,\partial_{\bar z},\qquad
 \partial=\partial_z,\qquad
 \partial_{\bar z}=\tfrac12(\partial_x+i\partial_y).
\]
The differential \(D\) has degree one, and \(\partial\) has degree zero.
They commute. Fields are smooth up to the boundary. Variations and test
fields have compact support, which may meet that boundary.

For partial forms of compact support use the traces
\begin{equation}\label{eq:vol2-native-traces}
 \tau_M(a)=\frac1{2\pi i}\int_M dz\wedge a,
 \qquad
 \tau_\partial(b)=\frac1{2\pi i}\int_{\mathbb C}dz\wedge b.
\end{equation}
They select partial form degree two and antiholomorphic degree one,
respectively. If \(i:\mathbb C\hookrightarrow M\) is the boundary inclusion,
these conventions give \(\tau_M(Da)=\tau_\partial(i^*a)\).
The sign includes both the outward normal and the position of \(dz\).

Let \(A=\mathbb C[u^a_n:1\leq a\leq r,\ n\geq0]\) be a differential
polynomial algebra, with \(\partial u^a_n=u^a_{n+1}\).
A finite-order matrix differential operator
\[
 H^{ab}(\partial)=\sum_{j=0}^N H_j^{ab}(u,\partial u,\ldots)\partial^j
\]
specifies generator brackets by
\(\{u^a_\lambda u^b\}=H^{ba}(\lambda)\).
Here \(\lambda\) records a holomorphic derivative of a contact.
The extension to all differential polynomials is
\begin{equation}\label{eq:vol2-native-master-formula}
 \{f_\lambda g\}=
 \sum_{a,b,m,n}
 \frac{\partial g}{\partial u^b_n}(\lambda+\partial)^n
 H^{ba}(\lambda+\partial)(-\lambda-\partial)^m
 \frac{\partial f}{\partial u^a_m}.
\end{equation}
Every operator acts on the expression to its right. The sum is finite on
each pair of polynomials.

We call \(H\) Hamiltonian when this bracket is skew and satisfies Jacobi.
Equivalently, \(H^{ab}=-(H^{ba})^*\), where the formal adjoint reverses
composition and sends \(\partial\) to \(-\partial\), and Jacobi holds on
generator triples. The master formula then propagates the identities to
all polynomials \cite[Theorem 1.15 and Proposition 1.16]{BarakatDeSoleKac2009}.
In particular it gives
\[
 \{\partial f_\lambda g\}=-\lambda\{f_\lambda g\},\qquad
 \{f_\lambda\partial g\}=(\lambda+\partial)\{f_\lambda g\}.
\]
The scalar operator \(H=\kappa\partial\) is the current example.

Introduce superfields \(\Phi^a\) and \(\eta_a\) of total degrees zero
and one. In their expansions in \(1,dt,d\bar z,dt\,d\bar z\), their
coordinate degrees are
\[
 \deg\Phi^a=(0,-1,-1,-2),\qquad
 \deg\eta_a=(1,0,0,-1).
\]
The degree-one zero-form in \(\eta_a\) is the gauge ghost. The negative
degrees are antifields. The BV two-form and action are
\begin{equation}\label{eq:vol2-native-action}
 \omega_{\mathrm{BV}}=\tau_M(\delta\Phi^a\delta\eta_a),
 \qquad
 S_H=\tau_M\left(\eta_aD\Phi^a+
              \tfrac12\eta_aH^{ab}(\partial;\Phi)\eta_b\right).
\end{equation}
The two-form has degree \(-1\), and the action has degree zero.
Its inverse pairing defines the degree-one BV bracket on local
functionals. This is the local Poisson vertex sigma model
\cite[\S\S2.2 and 4.1]{KhanZeng2025}.

\begin{proposition}\label{prop:vol2-native-classical}
For a Hamiltonian operator \(H\), the action
\eqref{eq:vol2-native-action} satisfies the classical master equation
on the boundary condition
\[
                         i^*\eta=0.
\]
This is a Lagrangian condition on the whole boundary superfield and is
preserved by the BV differential.
\end{proposition}

\begin{proof}
The kinetic square vanishes by \(D^2=0\). Its bracket with the interaction
is a total partial derivative. Its boundary trace contains \(i^*\eta\),
and vanishes. Pair the interaction bracket with three independent odd
test fields and integrate every holomorphic derivative by parts.
The remaining coefficient is the generator Jacobi expression
\[
 \{u^a_\lambda\{u^b_\mu u^c\}\}
 -\{u^b_\mu\{u^a_\lambda u^c\}\}
 -\{\{u^a_\lambda u^b\}_{\lambda+\mu}u^c\}.
\]
It vanishes by the Hamiltonian-operator identity. This proves
\(\{S_H,S_H\}_{\mathrm{BV}}=0\), and the graded Jacobi identity gives
\(Q^2=0\) for \(Q=\{S_H,-\}_{\mathrm{BV}}\).

Variation has boundary term
\(-\tau_\partial(\eta_a\delta\Phi^a)\).
The condition sets one complementary summand of this nondegenerate
pairing to zero and leaves the other free. It is Lagrangian.
The interaction part of \(Q\eta\) is quadratic in \(\eta\) and its
holomorphic derivatives. Its tangential pullback vanishes on the condition,
and \(\bar\partial\) preserves the same condition.
\end{proof}

For \(H=\kappa\partial\), write the physical fields as a scalar \(\phi\)
and a partial one-form \(a\). Then
\[
 S^{\mathrm{cl}}_\kappa=\tau_M\left(aD\phi+\tfrac\kappa2a\partial a\right),
 \qquad
 D\phi+\kappa\partial a=0,\qquad Da=0.
\]
The gauge variations are \(\delta\phi=\kappa\partial\epsilon\) and
\(\delta a=-D\epsilon\), with \(\epsilon\) zero at the boundary.
For \(\kappa\ne0\), the one-form
\(\mathcal A=a-\kappa^{-1}\phi\,dz\) identifies this action with
\(-\kappa(4\pi i)^{-1}\int\mathcal A\wedge d\mathcal A\).
Indeed expansion and one integration by parts give the displayed scalar
action, with boundary term zero under \(i^*a=0\). Thus the current model
has a pure abelian Chern--Simons realization with this normalization.

The finite Hamiltonian Lie algebra \(\mathfrak p_N\) of
Chapter~\ref{ch:vol2-hamiltonian-jets} supplies another classical example.
If \([e_a,e_b]=f_{ab}^{\ c}e_c\), set
\(\{u^a_\lambda u^b\}=f_{ab}^{\ c}u^c\).
This has differential order zero and satisfies Jacobi by the Lie identity.
Its action is the shifted cotangent action with those coefficients.
The finite-dimensional quantum measure of that chapter concerns its
algebraic model. Quantizing all spacetime modes and identifying a physical
boundary algebra remain additional comparisons.

\section{Regular Gaussian observables}

Keep \(\kappa\) as an indeterminate and put
\(R=\mathbb C[\kappa,\hbar]\). For the scalar Gaussian model use the
tangent complex
\[
 E=\left(\Omega_{\mathrm{HT}}^\bullet e_\Phi
       \oplus\Omega_{\mathrm{HT}}^\bullet e_\eta\right)\otimes\mathbb C[\kappa],
 \qquad |e_\Phi|=0,\quad |e_\eta|=-1.
\]
For a partial form \(\alpha\) of degree \(p\), its differential is
\begin{equation}\label{eq:vol2-native-tangent}
 \begin{split}
 Q(\alpha e_\Phi)&=D\alpha e_\Phi,\\
 Q(\alpha e_\eta)&=D\alpha e_\eta
            +(-1)^{p+1}\kappa\partial\alpha e_\Phi.
 \end{split}
\end{equation}
The two off-diagonal terms in \(Q^2\) have opposite signs. Impose the
relative boundary condition on every tangential \(e_\eta\) component.
The coordinate and tangent degrees are dual conventions.

The degree-minus-one pairing is
\[
 B(\alpha e_\Phi,\beta e_\eta)=\tau_M(\alpha\beta),\qquad
 B(\alpha e_\eta,\beta e_\Phi)
       =-(-1)^{\deg\beta}\tau_M(\alpha\beta),
\]
and it is zero on equal summands. Let \(E_c(U)\) denote the compactly
supported fields on an open \(U\), with the boundary condition when
needed. On \(C(U)=E_c(U)[1]\), write
\[
 d(s\xi)=-sQ\xi,\qquad
 b(s\xi,s\zeta)=(-1)^{|\xi|+1}B(\xi,\zeta).
\]
Thus \(b\) is a graded-symmetric form of degree one.

At fixed field number, a regular observable is a compact smooth joint
kernel in the corresponding graded symmetric tensor power of \(C(U)\).
Complete that tensor at fixed compact support in its smooth seminorms,
then take the inductive limit over compact supports. Denote the result
by \(\operatorname{Sym}_b^m C(U)\). Equivalently one can use the completed
bornological tensor of the compact smooth spaces. The direct sum over
field number remains algebraic. Define
\begin{equation}\label{eq:vol2-native-quantobs}
 \mathcal O^q(U)=
 \left(\bigoplus_{m\geq0}\operatorname{Sym}_b^m C(U)[\hbar],
                       d+\hbar\Delta_b\right).
\end{equation}
Here \(\Delta_b\) contracts an unordered pair with \(b\), with its Koszul
sign. Restriction of a smooth joint kernel to a diagonal and compact
integration define this operation continuously.

Compatibility of the pairing with \(Q\) gives
\([d,\Delta_b]=0\), and odd pair contractions give \(\Delta_b^2=0\).
Consequently \eqref{eq:vol2-native-quantobs} is a cochain complex.
Disjoint opens have products by extension by zero and concatenation of
kernels. There are no cross BV contractions on disjoint supports, so
these products commute with the differential and compose associatively.
The constant observable is their unit.

The descent statement uses the free-field theorem for an elliptic complex
of finite total rank, topological in the normal direction, with a local
Lagrangian elliptic boundary condition
\cite[Definitions 2.1--2.3 and Theorems 4.1--4.2]{GwilliamRabinovichWilliams2021}.
Our tangential differential is the triangular extension of two Dolbeault
complexes, and the normal factor is de Rham. These are precisely the
normally topological Gaussian data. The theorem gives a factorization
algebra of these regular observables and its boundary restriction.
The maps below construct the collar comparison explicitly.

Point evaluation has a delta-distribution kernel and is absent from
\eqref{eq:vol2-native-quantobs}. Its extension to local composite
observables is a renormalization problem. Likewise, the formal completion
\(\mathcal O^q(U)^{\wedge}_{\hbar}\) allows the field-number bound to grow
with the coefficient order. It is larger than completing each fixed
field number and then taking their algebraic sum.

\section{The normal contraction and its quantum correction}

Choose a smooth profile \(\rho\) of integral one supported in
\([0,\varepsilon)\), and put \(F(t)=\int_0^t\rho(s)\,ds\).
For a tangential open \(V\), split its compact fields into
\[
 L=\Omega_c^{0,\bullet}(V)e_\Phi\otimes\mathbb C[\kappa],
 \qquad G=\Omega_c^{0,\bullet}(V)e_\eta\otimes\mathbb C[\kappa].
\]
The tangential differential is
\(d_\partial(\ell,g)=(d_L\ell+Q_{\mathrm{rel}}g,d_Gg)\), where
\[
 Q_{\mathrm{rel}}(\alpha e_\eta)
       =(-1)^{\deg\alpha+1}\kappa\partial\alpha e_\Phi,
 \qquad d_LQ_{\mathrm{rel}}+Q_{\mathrm{rel}}d_G=0.
\]
Write a collar field as \(\gamma(t)+dt\,\beta(t)\), with
\(\gamma_G(0)=0\). Its differential is
\(d_\partial\gamma+dt(\partial_t\gamma-d_\partial\beta)\).

Define maps to and from \((G[-1],-d_G)\) by
\begin{equation}\label{eq:vol2-native-normal-maps}
 P(\gamma+dt\,\beta)=\int_0^\infty\beta_G(s)\,ds,
 \qquad Ig=dt\,\rho g+(F-1)Q_{\mathrm{rel}}g.
\end{equation}
The second term in \(I\) is required. Its normal derivative cancels the
unwanted term \(-dt\,\rho Q_{\mathrm{rel}}g\) in the differential of the
first term. In particular \(QI=-Id_G\) and \(PI=1\).

\begin{lemma}\label{lem:vol2-native-normal}
The degree-minus-one map
\[
 h(\gamma+dt\,\beta)=-\int_t^\infty\beta(s)\,ds
                  +(1-F(t))\int_0^\infty\beta_G(s)\,ds
\]
satisfies
\[
 Qh+hQ=1-IP,
 \qquad hI=Ph=h^2=0.
\]
These identities hold continuously on compact smooth joint kernels and
retain compact tangential support. No inverse of \(\kappa\) is used.
\end{lemma}

\begin{proof}
The correction sets the \(G\)-component of \(h\) to zero at \(t=0\).
All tail integrals vanish above the supports of the input and \(\rho\).
Differentiating the tail gives \(\beta(t)\), while differentiating
\(1-F\) gives \(-\rho\). Substitution in the collar differential proves
\(Qh+hQ=1-IP\). Since \(I\) has normal component \(\rho g\), its two
terms under \(h\) cancel. The output of \(h\) has no normal one-form,
so \(Ph=h^2=0\). Integrals in a compact normal interval are bounded in
each smooth tangential seminorm. The same estimates with other kernel
variables as parameters prove the completed-kernel statement.
\end{proof}

Let \(\mathfrak a(V)=\Omega_c^{0,\bullet}(V)\otimes\mathbb C[\kappa]\).
Fix \(\zeta^2=-1\). The map
\(\alpha\mapsto\zeta(-1)^q\alpha e_\eta\), for Dolbeault degree \(q\),
identifies \((\mathfrak a(V),\bar\partial)\) with \((G[-1],-d_G)\).
After suspension, \eqref{eq:vol2-native-normal-maps} and \(-sh\) give
a contraction \((i,p,h_C)\) from \(C\) to \(\mathfrak a(V)[1]\).

The pairing induced by \(i\) is determined by
\begin{equation}\label{eq:vol2-native-current-cocycle}
 \mu_\kappa(\alpha,\beta)
       =-\kappa\tau_\partial(\alpha\partial\beta),
 \qquad
 b(i(s\alpha),i(s\beta))=(-1)^q\mu_\kappa(\alpha,\beta).
\end{equation}
To check the normalization, only \(q+r=1\) contributes. The two terms
of \(I\) pair to
\(\kappa\tau_\partial(\alpha\partial\beta)\), since
\(2\int\rho(1-F)\,dt=1\) and
\(\tau_M(dt\,\omega)=-\int\tau_\partial(\omega)\,dt\).
The two phase factors give \(\zeta^2(-1)^{q+r}=1\), and suspension
gives \((-1)^{q+1}\). This proves the formula.

Write \(\mu^s=i^*b\) for the suspended pairing and set
\[
 \mathcal A_\kappa(V)=
 \left(\operatorname{Sym}_b(\mathfrak a(V)[1])[\hbar],
                  d+\hbar\Delta_{\mu^s}\right).
\]
Integration by parts makes \(\mu_\kappa\) a graded-skew cocycle.
The inclusion \(i^q=\operatorname{Sym}(i)\) therefore preserves the quantum
differential. The projection needs a correction.

\begin{theorem}\label{thm:vol2-native-quantum-collar}
Regular Gaussian collar observables contract onto
\(\mathcal A_\kappa(V)\) over \(R\). Both chain maps preserve the unit
and products on tangentially disjoint regions. They are continuous and
polynomial in \(\hbar\) at fixed field number, and extend coefficientwise
to the \(\hbar\)-adic completion.
\end{theorem}

\begin{proof}
Put \(e=ip\). Average \(h_C\otimes1+e\otimes h_C\) with its graded
transpose to obtain \(H_2\). Then
\(dH_2+H_2d=1-e\otimes e\). The even pairing \(P_2=-bH_2\) satisfies
\[
 \delta P_2=b-p^*\mu^s,
 \qquad [d,\partial_{P_2}]=\Delta_b-\Delta_{p^*\mu^s},
\]
where \(\partial_{P_2}\) contracts an unordered pair.
Constant contractions graded-commute. Hence for
\(W=\exp(\hbar\partial_{P_2})\) the conjugation has only one commutator:
\[
 W(d+\hbar\Delta_b)W^{-1}=d+\hbar\Delta_{p^*\mu^s}.
\]
Take
\begin{equation}\label{eq:vol2-native-quantum-maps}
 i^q=\operatorname{Sym}(i),\qquad
 p^q=\operatorname{Sym}(p)W.
\end{equation}
The pairing \(P_2\) is zero on \(i\otimes i\), so \(Wi^q=i^q\) and
\(p^qi^q=1\).

Split \(C=\operatorname{im}i\oplus\ker p\). Extend \(h_C\) as a derivation
of the symmetric algebra and divide by the positive number of \(\ker p\)
factors. Set the resulting \(H_0\) to zero on the retained factors alone.
Then \([d,H_0]=1-\operatorname{Sym}(e)\). Only retained factors enter
\(\Delta_{p^*\mu^s}\), so it graded-commutes with \(H_0\). Consequently
\[
 H^q=W^{-1}H_0W,
 \qquad (d+\hbar\Delta_b)H^q+H^q(d+\hbar\Delta_b)=1-i^qp^q.
\]
Each contraction lowers field number by two. All exponentials therefore
terminate on every polynomial. Normal integration and diagonal restriction
are continuous on smooth joint kernels. Their tangential supports remain
those of the uncontracted variables. In particular \(P_2\) has no cross
term on disjoint tangential supports, proving the product assertion.
At a fixed coefficient of \(\hbar\), only finitely many coefficients of a
formal input contribute. This proves the completed assertion as well.
\end{proof}

\section{Currents and supported interior insertions}

The finite-circle-spectrum sector on a centered disk is the algebraic
sum of its rotation eigenspaces. Choose a radial annular cutoff \(\chi\),
equal to one near the inner circle and zero beyond the outer circle.
The compact forms \(\alpha_n=(\bar\partial\chi)z^{-n}\), \(n\geq1\), obey
\[
 \tau_\partial(z^m\alpha_n)=\delta_{m,n-1},\qquad
 -\partial\alpha_n=n\alpha_{n+1}
                    -\bar\partial(z^{-n}\partial\chi).
\]
Stokes' formula proves the first identity. Differentiation proves the second.
They fix the creation modes and translation, including their normalization.

\begin{proposition}\label{prop:vol2-native-heisenberg}
The Gaussian boundary has polynomial vacuum states
\[
 V_\kappa=R[x_1,x_2,\ldots],\qquad
 J_{-n}=x_n,\quad J_n=\hbar\kappa n\partial_{x_n}\ (n>0),\quad J_0=0.
\]
Its current satisfies \(J(z)J(w)\sim\hbar\kappa(z-w)^{-2}\).
Translation is \(T=\sum_{n>0}n x_{n+1}\partial_{x_n}\).
These formulas include the nonsimple zero-level specialization.
\end{proposition}

\begin{proof}
The cocycle \eqref{eq:vol2-native-current-cocycle} on Laurent test functions
is \(-\kappa\operatorname{Res}(f\partial g\,dz)\). Thus the contour modes
obey \([J_m,J_n]=\hbar\kappa m\delta_{m+n,0}\).
The compact representatives supply the polynomial creation states and
\([T,J_m]=-mJ_{m-1}\). Normal-ordered products of derivatives of \(J\)
then give every state field, with the vacuum and translation identities.
Wick contraction proves locality and the stated double pole.

The action is quadratic. If its level term is treated as an insertion,
each kinetic propagator joins \(\Phi\) to \(\eta\), while that insertion
has two \(\eta\) legs. Two insertions cannot join each other.
A boundary pair receives exactly one insertion, and its two attachments
cancel the factor \(1/2\). Hence the current level has no interaction
correction. At \(\kappa=0\), the positive modes vanish and the same
polynomial carrier is a commutative differential vertex algebra.
\end{proof}

The current algebra does not select the Hilbert-space preparation of
Chapter~\ref{ch:vol2-infinite-mode-states}. Even in this Gaussian theory,
normalized polynomial correlations leave determinant and framing factors
undetermined. They are additional data of a physical state.

For an interior open \(U\subset V\times(a,b)\) with
\(0<\varepsilon<a<b\), the collar inclusion \(i^q\) has normal support
disjoint from \(U\). Define
\begin{equation}\label{eq:vol2-native-supported-action}
 \beta_{U;V}(f)(a_0)
       =p^q\bigl(\operatorname{ext}_U(f)i^q(a_0)\bigr).
\end{equation}
Its target is the complex of continuous operations on
\(\mathcal A_\kappa(V)\), with Hom differential
\(T\mapsto dT-(-1)^{|T|}Td\). Both comparison maps are cochain maps.
Disjointness removes the cross BV term in the middle product.
Thus \(\beta_{U;V}\) is a cochain map.

This action detects a contraction that linear restriction misses.
Take a unit-integral normal profile \(\rho_{\mathrm{int}}dt\) above the
boundary profile, and compact tangential tests \(\phi,\beta\) of
Dolbeault degrees zero and one with \(\tau_\partial(\phi\beta)=1\).
Set
\[
 u=\zeta^{-1}\kappa s(\rho_{\mathrm{int}}dt\,\phi e_\Phi),
 \qquad v=s\beta.
\]
Then \(p(u)=0\), but
\(h_Cu=\zeta^{-1}\kappa s(\phi e_\Phi)\) on the boundary profile.
Substitution in \(P_2=-bH_2\) gives
\[
 P_2(u,i(v))=\kappa,
 \qquad \beta(u)(v)=\hbar\kappa,
 \qquad \beta(u)(1)=0.
\]
The cochain \(u\) is not assumed closed. This is an operation on the
supported complexes, rather than a new bulk cohomology class.

For a homogeneous linear input, the complete formula is
\[
 \begin{split}
 \beta(u)(v_1\cdots v_r)
 &=p(u)v_1\cdots v_r\\
 &\quad+\hbar\sum_{j=1}^r
 (-1)^{|u|\sum_{\ell<j}|v_\ell|}P_2(u,i(v_j))
                  v_1\cdots\widehat v_j\cdots v_r.
 \end{split}
\]
Its quantum factor remains over every truncated or completed coefficient
ring. There is no division by \(\hbar\).

If \(f\) is supported above \(g\), the normal homotopy of \(g\) remains
below \(f\). Inserting \(1-i^qp^q=dH^q+H^qd\) gives the ordered
composition homotopy
\begin{equation}\label{eq:vol2-native-action-homotopy}
 K(f,g,a_0)=(-1)^{|f|}p^q\bigl(
       \operatorname{ext}(f)H^q(\operatorname{ext}(g)i^q(a_0))\bigr).
\end{equation}
Its Hom differential is the joint projected product minus the successive
actions. The sign moves the degree-minus-one homotopy past \(f\).
Arbitrary motion of supports requires parameter-dependent operations.
The forest coefficient homotopies of
Chapter~\ref{ch:vol2-supported-compositions} retain their contact terms,
but their decoration by native field operations is a further comparison.

\section{Smooth periods and polynomial states}

A smooth compact smearing need not have finitely many holomorphic moments.
For \(0<r<R_0\), define
\[
 E_r=\left\{\sum_{n\geq1}a_nw^{-n}:
                       \sum_{n\geq1}|a_n|r^{-n}<\infty\right\},
 \qquad \mathscr P(D_{R_0})=\operatorname*{colim}_{r<R_0}E_r.
\]
The norm is the displayed weighted sum. For \(r<r'\), the inclusion has
diagonal coefficients \((r/r')^n\) in normalized bases and is nuclear
and compact. A countable cofinal sequence of radii defines the complete
locally convex inductive limit. Its elements are exterior holomorphic
germs vanishing at infinity. Cauchy estimates identify this description
with the coefficient description above.

Put
\[
 \mathscr F(D_{R_0})=
       \bigoplus_{m\geq0}\operatorname{Sym}_\pi^m\mathscr P(D_{R_0})[\hbar].
\]
Here the projective tensor is completed at each fixed \(m\), and the
field-number sum is algebraic. A finite union of disjoint disks uses the
direct sum of their germ spaces. Common inner radii and compact support
bounds are required locally on parameter families.

The exterior Cauchy map on a compact \((0,1)\)-form is
\begin{equation}\label{eq:vol2-native-cauchy}
 \mathcal C(sa)(w)=\tau_\partial\left(\frac{a(z)}{w-z}\right)
             =\sum_{n\geq0}\tau_\partial(z^na)w^{-n-1}.
\end{equation}
It is zero on shifted functions. The expansion is taken outside the
support disk. A radial cutoff \(\chi\) provides a section
\(j_\chi(f)=s(\bar\partial\chi)f\), on any germ stage lying inside its
transition annulus.

\begin{lemma}\label{lem:vol2-native-dolbeault}
At every fixed field number, the symmetric Cauchy map is a continuous
quasi-isomorphism from the compact Dolbeault current complex to the
corresponding symmetric germ space. Its cohomology is in degree zero,
with its Hausdorff quotient topology, and its differential has closed
range.
\end{lemma}

\begin{proof}
The Cauchy convolution on the whole plane is a smooth function
\(\widetilde{\mathcal C}a\) with
\(\bar\partial\widetilde{\mathcal C}a=-a\).
If its exterior germ is zero, it is a compact primitive. There are no
nonzero compact holomorphic functions.
Choose radii
\(r_0<r<r'<r''<R_0\), with the source supported in \(D_{r_0}\) and
\(\chi=1\) through \(D_{r'}\), supported in \(D_{r''}\). Then
\[
 h_\chi(sa)=s\left(\widetilde{\mathcal C}a
                           -(1-\chi)\mathcal C(a)\right),
 \qquad h_\chi(s\phi)=0
\]
satisfies \(d_1h_\chi+h_\chi d_1=1-j_\chi\mathcal C\), where
\(d_1(s\alpha)=-s\bar\partial\alpha\).
The second term is defined outside \(D_r\) and extended by zero where
\(1-\chi=0\). Stokes' formula gives \(\mathcal Cj_\chi=1\).

Cauchy estimates on the annular gap bound every derivative by finitely
many source seminorms and the \(E_r\) norm. Apply this contraction in each
variable of a joint kernel. The tensor homotopy applies \(h_\chi\) to one
variable, the projection to earlier variables and the identity to later
variables, with Koszul signs. The estimates are uniform in the remaining
smooth parameters. Thus the identity holds on completed projective
tensors. Averaging over the finite permutation group gives symmetric
powers. The ranges are kernels of continuous maps to Hausdorff spaces,
and the stagewise sections make the inverse on cohomology continuous.
\end{proof}

The Gaussian and source terms can also be removed explicitly.
Let \(k_1(z)=1/z\) as a locally integrable distribution and
\(k_2=-\partial_zk_1\). On shifted top forms define
\[
 P_{\mathrm C}(sa,sb)=\kappa\tau_z\tau_w
                      \bigl(k_2(z-w)a(z)b(w)\bigr),
 \qquad \gamma(sa)=\tau_z\left(\frac{a(z)}z\right).
\]
They vanish on other form degrees. Their Hom differentials are
\(\delta P_{\mathrm C}=\mu^s\) and \(\delta\gamma=\epsilon_0\), where
\(\epsilon_0(s\phi)=\phi(0)\) and \(\epsilon_0\) is zero on top forms.
Indeed \(\partial_{\bar z}(1/z)=\pi\delta_0\) and
\(\tau(f\,d\bar z)=-\pi^{-1}\int f\) give
\[
 \tau\left(\frac{\bar\partial\phi}{z}\right)=\phi(0),
 \qquad
 \tau_z\bigl(k_2(z-w)\bar\partial\phi(z)\bigr)=\partial_w\phi(w).
\]

For a scalar Wilson charge \(q\), the boundary differential is
\(d_1+\hbar\Delta_{\mu^s}+\hbar q\epsilon_0\).
The exponential
\[
 N_q=\exp\bigl(\hbar\partial_{P_{\mathrm C}}
                       +\hbar q\partial_\gamma\bigr)
\]
conjugates it to \(d_1\). The proof is the same one-commutator calculation
as in Theorem~\ref{thm:vol2-native-quantum-collar}.
Every contraction removes factors. Products of contracted distributions
use disjoint pairs of variables, and therefore exist as distributional
tensor products. They are paired with smooth joint kernels.

The line source on bulk observables is the constant derivation
\[
 \ell(s\xi)=\zeta^{-1}\int_0^\infty\beta_\eta^{(0)}(t,0)\,dt,
 \qquad \xi=\gamma+dt\,\beta.
\]
Its degree is one. The relative boundary condition gives \(\ell d=0\),
and it graded-commutes with constant contractions. Thus
\(d+\hbar\Delta_b+\hbar q\ell\) squares to zero.
Moreover \(\ell=\epsilon_0p\), and the quantum collar maps still give a
contraction onto the charged boundary complex. The Gaussian identity
\(e^{-q\lambda}(d+\hbar\Delta_b)e^{q\lambda}
=d+\hbar\Delta_b+\hbar q\ell\) expresses the Wilson insertion, where
\(\lambda\) is its distributional linear representative. The identity
defines the operation on regular observables without placing \(\lambda\)
inside the smooth observable algebra.

Consequently
\begin{equation}\label{eq:vol2-native-smooth-period}
             \Pi=\operatorname{Sym}(\mathcal C)N_qp^q
\end{equation}
is a continuous quasi-isomorphism onto \(\mathscr F(D_{R_0})\), with zero
target differential. The statement follows by composing the three
explicit contractions. It is a comparison with the germ completion.
For example \(e^{c/w}-1\), for \(c\ne0\), is the Cauchy transform of
\(s(\bar\partial\chi)(e^{c/z}-1)\) and has infinitely many moments.

On a centered disk, rotation gives \(w^{-n}\) weight \(-n\).
At a fixed total negative weight there are finitely many partitions into
positive mode numbers. Hence the finite-circle-spectrum target is exactly
\(R[x_1,x_2,\ldots]\), with \(x_n=w^{-n}\).
The restricted period map remains a quasi-isomorphism. To see this, choose
a primitive of a finite-spectrum cocycle in its full acyclic cone, then
apply the finite Fourier projector onto the weights of that cocycle.
All rotations lie in a common compact support and germ stage, so the
integral exists and commutes with the differential. This gives a
finite-spectrum primitive.

The same map identifies separated current insertions with
\[
 J(z)=\sum_{n\geq1}x_nz^{n-1}+\frac{\hbar q}{z}
                  +\hbar\kappa\sum_{n\geq1}n z^{-n-1}\partial_{x_n}.
\]
A compact radial current smearing centered at \(z\), disjoint from the
inner support, has Cauchy transform \((w-z)^{-1}\). Its charge contraction
is \(\hbar q/z\), and its covariance with an inner period \(f\) is
\(-\hbar\kappa f'(z)\). These three terms give the formula.
Contour integration extracts each mode. Changing the contour radius within
the separated annulus is a cochain homotopy by Stokes' formula.

\section{An incoming junction from the Gauss constraint}

We now set \(\hbar=\kappa=1\) and retain a numerical charge \(q\).
Consider the incoming Wilson ray in complex abelian Chern--Simons theory
\[
 S=-\frac1{4\pi i}\int\mathcal A\wedge d\mathcal A,
 \qquad W_q=\exp\left(q\int_\gamma\mathcal A\right).
\]
On an oriented boundary disk with coordinate \(w\), orient the incoming
line opposite to the outward normal. Variation of the normal multiplier
then imposes
\begin{equation}\label{eq:vol2-native-gauss}
 d_Da=-2\pi iq\,\delta_0,
 \qquad a=a_q+\alpha,
 \qquad a_q=-q\,\frac{dw}{w},
\end{equation}
where \(\alpha\) is smooth and \(\delta_0\) has integral one.
The singularity class is fixed in this statement. Gauge functions vanish
on the boundary circle. There is no quotient by large gauge
transformations and no added monopole sector.

The fluctuation complex is
\[
 \Omega^0_\partial(D)\xrightarrow{d}\Omega^1(D)
                               \xrightarrow{d}\Omega^2(D)
\]
in degrees \(-1,0,1\). It retracts to
\(C^\infty(S^1)/\mathbb C\) in degree zero. One explicit contraction uses
radial Poincar\'e maps \(h_1,h_2\), with
\(dh_1+h_2d=1\), \(dh_2=1\), \(h_1h_2=0\), and an extension \(E\) of
mean-zero boundary functions. Put
\[
 P\alpha=[\operatorname{tr}h_1\alpha],\qquad Ib=dEb,
 \qquad
 H_0\alpha=h_1\alpha-E(P\alpha)
                    -\operatorname{mean}(\operatorname{tr}h_1\alpha),
 \quad H_1=h_2.
\]
The Poincar\'e identities give \(PI=1\) and \(dH+Hd=1-IP\).
For a gauge function vanishing at the circle, the mean term removes the
constant in \(h_1df=f-f(0)\). Thus the contraction respects the relative
domain, not only the unrestricted disk complex.

The reduced form and currents are
\begin{equation}\label{eq:vol2-native-reduced-form}
 \Omega_{\mathrm{red}}
      =\sum_{n>0}\frac{\delta J_{-n}\wedge\delta J_n}{n},
 \qquad J_n=-\frac1{2\pi i}\oint w^na,
 \qquad J_0=q.
\end{equation}
Indeed the smooth solutions are \(a_q+df\), modulo zero-boundary gauge
functions. Stokes' formula gives
\(\Omega(\delta_1f,\delta_2f)
=(2\pi i)^{-1}\oint\delta_1f\,d\delta_2f\), which gives the nonzero
modes in \eqref{eq:vol2-native-reduced-form}. The source fixes the zero mode.

To retain the constraints before reduction, choose Dirichlet functions
\(\epsilon_j\) orthonormal for the energy pairing and put
\[
 u_j=d\epsilon_j,\qquad v_j=2\pi i\,*d\epsilon_j.
\]
For \(\Omega_D(\alpha,\beta)=(2\pi i)^{-1}\int_D\alpha\wedge\beta\),
these obey \(\Omega_D(u_j,v_k)=\delta_{jk}\) and zero equal-type pairings.
Harmonic gradients are symplectically orthogonal to these modes by
Green's formula. Solving the two Dirichlet Poisson equations gives the
exact, coexact and harmonic splitting.

Write the constraint part of \(\alpha\) as \(\sum_j(s_ju_j+p_jv_j)\).
Its moment is
\[
 -\frac1{2\pi i}\int_D\epsilon_j(d_Da+2\pi iq\delta_0)=p_j.
\]
The constraint \(p_j=0\) translates \(s_j\).
Keep finitely many gauge and boundary modes, quantize in the incoming
polarization \(x_n=J_{-n},s_j\), and take the union of the resulting
polynomial state spaces. With odd ghosts \(\nu_j\) of degree one this
defines
\begin{equation}\label{eq:vol2-native-bfv-junction}
 \mathcal J_q=
 \left(\mathbb C[x_1,x_2,\ldots,s_1,s_2,\ldots]
          \otimes\Lambda(\nu_1,\nu_2,\ldots),
       \quad Q_{\mathrm{BFV}}=\sum_j\nu_j\partial_{s_j}\right).
\end{equation}
All polynomials are finite. The differential is the quantized commuting
moment constraint \(\sum_j\nu_jp_j\). No constant gauge ghost occurs for
the chosen zero-boundary gauge group.

\begin{theorem}\label{thm:vol2-native-polynomial-junction}
Evaluation at zero gauge coordinates and ghosts gives a quasi-isomorphism
\[
 \rho_q:\mathcal J_q\longrightarrow F_q=\mathbb C[x_1,x_2,\ldots],
 \qquad \rho_q(f)=f(x,0,0).
\]
Its inverse up to homotopy is the inclusion \(\iota_q(f)=f(x)\).
These maps and their contraction commute with
\[
 J_{-n}=x_n,\qquad J_n=n\partial_{x_n},\qquad J_0=q,
 \qquad T_q=\sum_{n>0}n x_{n+1}\partial_{x_n}+qx_1.
\]
Thus they preserve every current field, the rooted residues and the
identity action of the local vacuum.
\end{theorem}

\begin{proof}
Let \(N\) count total polynomial and exterior degree in \(s,\nu\),
excluding \(x\). With left odd derivatives set
\[
 H=\frac1N\sum_js_j\partial_{\nu_j}\quad(N>0),\qquad H=0\quad(N=0).
\]
The anticommutator of its numerator with \(Q_{\mathrm{BFV}}\) is the
gauge Euler derivation. Therefore
\(Q_{\mathrm{BFV}}H+HQ_{\mathrm{BFV}}=1-\iota_q\rho_q\).
The sum is finite on each state and passes to the algebraic union of
cutoffs. Every displayed current and \(T_q\) acts only on \(x\), so it
commutes with the contraction. Normal ordering and differentiation give
all current fields. Their Laurent coefficients, including coefficients
tested by multipliers before taking principal parts, give the rooted
residues. The vacuum acts as the identity throughout.
\end{proof}

For the unnormalized disk function \(\epsilon=1-r^2\), its energy is
\(2\pi\). Thus \(u=d\epsilon\), \(v=i*d\epsilon\) form one canonical
pair and give the differential \(\nu\partial_s\).
Removing the Wilson source leaves this contractible gauge complex
unchanged but sets \(J_0=0\). It then has zero-mode defect two on the
unit of the charge-two recipient. Gauge contraction alone does not
recover the missing charge.

\section{The charged operation and its collision lattice}

Place a second unit charge at separation \(z\ne0\).
The source in \eqref{eq:vol2-native-gauss} changes to
\(-dw/w-dw/(w-z)\). Contour integration gives the current Ward identity
\[
 J_m^{\mathrm{out}}\mathcal V(z)-\mathcal V(z)J_m^{\mathrm{in}}
       =z^m\mathcal V(z),\qquad m\in\mathbb Z.
\]
In the incoming polarization, its oscillator solution with highest
coefficient one is
\begin{equation}\label{eq:vol2-native-endpoint-kernel}
 K(z)f(x)=\exp\left(\sum_{n>0}\frac{x_nz^n}{n}\right)
                                      f(x_n-z^{-n}).
\end{equation}
The creation exponential supplies the positive-mode Ward term, and the
shift supplies the negative-mode term. On a polynomial input, the shift
is finite. At a fixed output weight, only finitely many creation modes
contribute. Values lie in the product of oscillator-weight spaces.

Translation fixes the remaining scalar. Direct substitution gives
\[
 T_2K-KT_1=\partial_zK+z^{-1}K.
\]
For a charged operation \(h(z)K(z)\), translation covariance requires
\(h'=h/z\). Normal ordering removes each endpoint's self-contraction,
and normalization at \(z=1\) gives
\begin{equation}\label{eq:vol2-native-charged-operation}
                         \mathcal V(z)=zK(z).
\end{equation}
Insertion of negative current modes gives every descendant of the first
input. Explicitly, if \(I(u,z)\) denotes an endpoint operation, then
\[
 I(J_{-n}u,z)=
 \frac{\partial_z^{n-1}J^-(z)}{(n-1)!}I(u,z)
 +I(u,z)\frac{\partial_z^{n-1}J^+(z)}{(n-1)!}.
\]
Here \(J^-\) contains negative modes, while \(J^+\) contains nonnegative
modes. The indicated left and right operations commute for different
\(n\), so this gives a well-defined extension to every polynomial state.
Its current commutator follows from
\(m\binom{m-1}{n-1}=n\binom mn\). Thus the operation retains descendants,
not just the highest coefficient.

Let \(\mathcal L^-\) be the operation line with normalized frame \(k=K\).
The connection on operations is
\(\nabla^{\mathrm{op}}I=dI-T_2I\,dz+IT_1\,dz\), and hence
\[
 \nabla^-k=-\frac{dz}{z}k,
 \qquad \mathcal L^+=(\mathcal L^-)^\vee,
 \qquad \nabla^+k^\vee=\frac{dz}{z}k^\vee.
\]
The fused state recipient is \(F_2\otimes\mathcal L^+\).
To verify its representing property on finite sums of scalar-charge Fock
modules, apply the zero-mode Ward identity first: the output charge must
be two. Removing the creation exponential and shift leaves a current-linear
map of polynomial Fock modules. Such a map is determined by a scalar on
the highest vector. The descendant recursion supplies every first-input
state. The resulting operation space is
\(\mathcal L^-\otimes\operatorname{Hom}_{\mathrm{cur}}(F_2,W)\), which
is the Hom space from the stated dual-line recipient, with its Hom
connection.

The map \(\iota_2\) therefore gives a quasi-isomorphism of flat families
\[
 F_2\otimes\mathcal L^+
       \longrightarrow \mathcal J_2\otimes\mathcal L^+.
\]
The physical Gaussian operation is carried to \(\mathcal V\) by
\(\rho_2\). Counterclockwise half-transport is \(-1\) on the state line
and on its dual operation line. For general scalar charges \(a,b\) at
level \(\kappa\), the same calculation gives opposite connection forms
\(\pm(ab/\kappa)d\log z\). Their transports are dual, and need not be
equal when \(ab/\kappa\) is nonintegral.

There is a specific extension through collision for the recipient.
Extend the normalized frames over \(\mathbb C[z]\), defining
\[
 \overline{\mathcal E}_F=\mathbb C[z]\otimes F_2\,k^\vee,
 \qquad
 \overline{\mathcal E}_{\mathrm{phys}}
       =\mathbb C[z]\otimes\mathcal J_2\,k^\vee.
\]
The connection is logarithmic:
\(\nabla_{z\partial_z}(a(z)m)=(za'(z)+a(z))m\).
The contraction is independent of \(z\), so
\(1\otimes\iota_2\) is a quasi-isomorphism over \(\mathbb C[z]\),
preserves this connection and has residue \(+1\). Its derived fibre at
zero is the quasi-isomorphism \(F_2\to\mathcal J_2\).
The horizontal frame \(z^{-1}k^\vee\) specifies a different lattice with
residue zero. A meromorphic change of frame cannot be evaluated at zero
to identify their chosen fibres.

The operation has a different limit. If \(f\) has oscillator weight \(w\),
then its output of weight \(d\) has the form
\[
 \operatorname{pr}_d\mathcal V(z)f=z^{1-w+d}v_d(f).
\]
An annihilation removes weight \(a\) with factor \(z^{-a}\), a creation
adds weight \(c\) with factor \(z^c\), and \(d=w-a+c\).
Consequently, separately in every output weight,
\[
 \lim_{z\to0}z^{w-1}\mathcal V(z)f
                     =f(-1,-1,\ldots)|2\rangle.
\]
There is no common finite-power normalization for all inputs:
\begin{equation}\label{eq:vol2-native-unbounded-input-poles}
       \langle2|\mathcal V(z)x_1^m=(-1)^mz^{1-m}.
\end{equation}
The first three values are \(z,-1,z^{-1}\).
Removing the negative shift would remove these poles and violate the
\(J_{-1}\) Ward identity. The framed recipient fibre and the
weight-dependent operation limit are therefore distinct constructions.

\section{Regular Wilson jets and corrected cutoff maps}

Return to a nonzero numerical \(\kappa\), keeping \(\hbar=1\).
For the incoming face use the pairing \(b_-=-b\) and the linear inclusion
\(i_- =\zeta^{-1}i\). These changes preserve the boundary pairing:
\(b_-(i_-\otimes i_-)=\mu^s\).
The ray functional then has no phase factor.
For \(\xi=\gamma+dt\,\beta\), define constant derivations
\begin{equation}\label{eq:vol2-native-ray-jets}
 \begin{aligned}
 \ell_n(s\xi)&=\frac1{n!}\int_0^\infty
                  \partial_z^n\beta_\eta^{(0)}(t,0)\,dt,& |\ell_n|&=1,\\
 r_n(s\xi)&=\frac1{n!}\int_0^\infty
                  \partial_z^n\beta_\Phi^{(0)}(t,0)\,dt,& |r_n|&=0,\\
 e_n(s\xi)&=\frac1{n!}\partial_z^n\gamma_\Phi^{(0)}(0,0),& |e_n|&=1.
 \end{aligned}
\end{equation}
The superscript selects Dolbeault degree zero. Each derivation evaluates
one variable of a smooth kernel. No product of point distributions is
introduced.

\begin{lemma}\label{lem:vol2-native-endpoint}
The ray and endpoint derivations satisfy
\begin{equation}\label{eq:vol2-native-endpoint-transgression}
 [d,\ell_n]=[d,e_n]=0,
 \qquad [d,r_n]=\kappa(n+1)\ell_{n+1}-e_n.
\end{equation}
They graded-commute with each other and with \(\Delta_{b_-}\).
\end{lemma}

\begin{proof}
On a linear test field, suspension and normal integration give
\[
 r_n(d(s\xi))=e_n(s\xi)-\kappa(n+1)\ell_{n+1}(s\xi).
\]
The endpoint term uses
\(\int_0^\infty\partial_t\gamma_\Phi^{(0)}dt
=-\gamma_\Phi^{(0)}(0)\). The other term is the off-diagonal part of
\eqref{eq:vol2-native-tangent}, including the derivative factorial.
Since \(r_n\) has degree zero, its commutator with \(d\) is \(-r_nd\)
on a linear input. The relative condition kills the possible endpoint
terms for \(\ell_n d\) and \(e_nd\). Constant contractions graded-commute.
Continuity extends these calculations to smooth joint kernels.
\end{proof}

For a finite set \(F\subset\mathbb Z_{>0}\), introduce the regular
coefficient algebra
\[
 Y_F=\mathbb C[B_0,B_n:n\in F]\otimes\Lambda(c_{n-1}:n\in F),
 \quad Q_YB_n=\kappa n c_{n-1},\quad Q_YB_0=Q_Yc_j=0.
\]
The \(B\)'s have degree zero and the \(c\)'s degree one. Tensors with
\(Y_F\) are algebraic in the coefficients. Operator expressions use
graded left multiplication, while the retained module structure is the
regular right action.
Set
\[
 D_0=d+\Delta_{b_-}+Q_Y+B_0\ell_0,
 \qquad L_F=\sum_{n\in F}\frac{B_n}{\kappa n}r_{n-1}.
\]
The bare descendant source would add
\(\sum_{n\in F}(B_n\ell_n+c_{n-1}r_{n-1})\) to \(D_0\).
Its square is \(\sum_{n\in F}c_{n-1}e_{n-1}\).
For the first jet it takes the value \(c_0\) on a suspended test function
equal to one at the endpoint. This is a nonzero boundary defect.

Define the corrected differential by
\begin{equation}\label{eq:vol2-native-corrected-source}
 D_F=D_0+\sum_{n\in F}
       \left(B_n\ell_n+c_{n-1}r_{n-1}
                         -\frac{B_n}{\kappa n}e_{n-1}\right).
\end{equation}
The correction is forced by the endpoint transgression.

\begin{proposition}\label{prop:vol2-native-source-gauge}
On regular smooth observables with coefficients in \(Y_F\),
\[
                 D_F=e^{-L_F}D_0e^{L_F},\qquad D_F^2=0.
\]
The conjugation is right \(Y_F\)-linear, fixes the unit and is finite
on every input of finite field number.
\end{proposition}

\begin{proof}
The term \(Q_YB_n=\kappa nc_{n-1}\) cancels the interior part of
\([d,c_{n-1}r_{n-1}]\). Its endpoint part is the bare defect above.
Differentiating \(-B_ne_{n-1}/(\kappa n)\) cancels that defect.
More precisely,
\[
 [D_0,L_F]=D_F-D_0,
 \qquad [L_F,[D_0,L_F]]=0.
\]
The second identity follows from constant contractions and commuting
even coefficients. The conjugation formula follows. Each \(r_n\)
lowers field number by one and kills constants, so its exponential
terminates and fixes \(1\otimes Y_F\).
\end{proof}

Let \((i_0,p_0,H_0)\) be the scalar-source quantum collar contraction in
the incoming phase, extended in the remaining coefficients. Then
\[
 I_F=i_0,
 \qquad P_F=p_0e^{L_F},
 \qquad H_F=e^{-L_F}H_0e^{L_F}
\]
contract \((\mathcal O(U)\otimes Y_F,D_F)\) onto
\[
 \mathcal A_F(V)=
 \left(\operatorname{Sym}_b(\Omega_c^{0,\bullet}(V)[1])\otimes Y_F,
         d_1+\Delta_{\mu^s}+B_0\epsilon_0+Q_Y\right).
\]
Indeed \(L_Fi_0=0\): the normal inclusion has no normal
\(e_\Phi\)-component, and its Wick contractions commute with \(L_F\).
Conjugation now transports every strong contraction identity.
The maps preserve the unit, extension of compact tangential supports
and products on tangentially disjoint regions. In particular all
\(B_n,c_j\) continue to act regularly.

Adding a jet changes the differential of a raw smooth observable.
For \(F\subset G\), coefficient inclusion \(\iota_{F,G}\) is therefore
replaced by
\begin{equation}\label{eq:vol2-native-cutoff-transition}
 j_{F,G}=e^{-(L_G-L_F)}\iota_{F,G}.
\end{equation}
The commuting \(L\)'s give
\(D_Gj_{F,G}=j_{F,G}D_F\) and \(j_{G,H}j_{F,G}=j_{F,H}\).
The maps are injective. With \(Y_\infty=\operatorname{colim}_F Y_F\),
their direct limit has the explicit untwisting
\begin{equation}\label{eq:vol2-native-all-jet-source}
 \Psi:\operatorname{colim}_F(\mathcal C_F,j_{F,G})
       \xrightarrow{\sim}(\mathcal O(U)\otimes Y_\infty,D_0),
 \qquad [f]_F\longmapsto e^{L_F}f.
\end{equation}
For its inverse, choose a finite stage containing an input's coefficient
indices and send \(a\) to \([e^{-L_F}a]_F\). Enlarging the stage gives
the same class. This proves the isomorphism and defines the action of
every coefficient. It is continuous for the final topology of finite
coefficient, field-number and compact smooth stages.

The normal contractions commute with these transitions, since
\(P_Gj_{F,G}=P_F\), \(j_{F,G}I_F=I_G\) and
\(H_Gj_{F,G}=j_{F,G}H_F\) after coefficient extension.
They consequently retract the direct limit onto \(\mathcal A_\infty\).
No inverse-limit cohomology assertion is used.

This construction does not apply the literal infinite source to an
unchanged smooth kernel. On a disk of radius less than one choose
\(\chi=1\) near zero and an interior normal profile of integral one.
The linear observable
\[
 u=s\left(dt\,\rho(t)\frac{\chi(z)}{1-z}e_\Phi\right)
\]
has \(r_n(u)=1\) for every \(n\), and \(\ell_n(u)=e_n(u)=0\).
For \(F=\{1,\ldots,K\}\),
\[
 D_Fu=du+\sum_{j=0}^{K-1}c_j,
 \qquad e^{L_F}u=u+\sum_{n=1}^K\frac{B_n}{\kappa n}.
\]
Neither infinite sum belongs to \(Y_\infty\).
The representative of \([u]_\varnothing\) at stage \(F\) is instead
\(u-\sum_{n\in F}B_n/(\kappa n)\).
Quotient cutoff maps also define an inverse-limit carrier, but that
carrier admits coefficient series absent from the direct limit above.

\section{The incoming cap and its state domain}

Geometric rotations act on the incoming test fields by
\[
 R_\theta(s\alpha e_\eta)=s\rho_\theta^*\alpha e_\eta,
 \qquad
 R_\theta(s\alpha e_\Phi)=e^{i\theta}s\rho_\theta^*\alpha e_\Phi,
 \quad \rho_\theta(t,z)=(t,e^{i\theta}z),
\]
and give \(B_n,c_{n-1}\) weight \(n\), with \(B_0\) of weight zero.
Pullback includes the form factors. The extra factor on \(e_\Phi\)
makes the off-diagonal \(\partial_z\) equivariant. Direct evaluation
shows that every source term, \(L_F\), cutoff map and normal contraction
commutes with this circle action.

Define \(\mathcal J_{\mathrm{rot}}(D_{R_0})\) as the algebraic sum of
the circle eigenspaces in the direct-limit source. This is a geometric
condition before taking periods. Combining
\eqref{eq:vol2-native-smooth-period} with \(\Psi\) gives a continuous
right \(Y_\infty\)-linear quasi-isomorphism
\begin{equation}\label{eq:vol2-native-regular-period}
 \Pi_{\mathrm{rot}}:\mathcal J_{\mathrm{rot}}(D_{R_0})
       \longrightarrow
       (\mathbb C[x_1,x_2,\ldots]\otimes Y_\infty,Q_Y).
\end{equation}
For clarity, the period at stage \(F\) is
\(\operatorname{Sym}(\mathcal C)
\exp(\partial_{P_{\mathrm C}}+B_0\partial_\gamma)p_0e^{L_F}\).
All exponentials are finite on each source input.

To prove the assertion about finite circle spectrum, each coefficient
element uses finitely many monomials of \(Y_\infty\).
Separating their weights reduces the target to the finite-weight germ
calculation above. The finite Fourier projector argument then restricts
the full quasi-isomorphism. Separated current insertions give
\[
 J_{-n}=x_n,\qquad J_n=\kappa n\partial_{x_n},\qquad J_0=B_0.
\]
The same Cauchy, covariance and charge contractions prove this on
cochains, before passing to cohomology.

The state produced by an incoming cap is a further object.
For fixed finite \(F\), write the boundary superfields as
\(\Phi_\Sigma=\phi+d\bar z\,\phi_b\) and
\(\eta_\Sigma=c+d\bar z\,a_b\). The boundary jet is
\(x_n=\partial_z^{n-1}\phi(0)/(n-1)!\).
The incoming boundary charge, with its endpoint, is
\begin{equation}\label{eq:vol2-native-incoming-charge}
 \begin{split}
 S^F_{\mathrm{in}}
 &=-\frac1\pi\int c(\partial_{\bar z}\phi+\kappa\partial_z a_b)\,dx\,dy
       +B_0c(0)\\
 &\quad+\sum_{n\in F}
       \left(\frac{B_n\partial_z^nc(0)}{n!}+c_{n-1}x_n\right).
 \end{split}
\end{equation}
The line coefficients are Casimirs for the boundary pairing. Their
differential \(Q_Y\) is retained externally. The incoming form and
potential have the opposite signs to those of the outward face.
In the polarization with \(\eta_\Sigma\) as momentum directions,
\[
             \widehat a_b=\pi\frac\delta{\delta\phi},
             \qquad \widehat c=\pi\frac\delta{\delta\phi_b}.
\]
Derivatives are placed to the right.

The relative master identity can be checked on smeared fields. For a
test function \(g\), the induced derivation has
\[
 \begin{aligned}
 q_F\phi(g)&=-\kappa c(\partial_zg),&q_Fc(g)&=0,\\
 q_Fa_b(g)&=c(\partial_{\bar z}g)
       -\pi\sum_{n\in F}c_{n-1}
                         \frac{\partial_z^{n-1}g(0)}{(n-1)!},\\
 q_F\phi_b(g)&=\phi(\partial_{\bar z}g)+\kappa a_b(\partial_zg)
       +\pi B_0g(0)+\pi\sum_{n\in F}B_n\frac{\partial_z^ng(0)}{n!}.
 \end{aligned}
\]
Together with \(q_F|_{Y_F}=Q_Y\), its square vanishes.
On \(\phi_b(g)\), the two bulk derivatives cancel,
and the endpoint cancellation is
\[
 -\pi\kappa\sum_n c_{n-1}\frac{\partial_z^ng(0)}{(n-1)!}
 +\pi\sum_n\kappa nc_{n-1}\frac{\partial_z^ng(0)}{n!}=0.
\]
Omitting the line-ghost term \(c_{n-1}x_n\) leaves the second sum.

Let \(\mathcal P_{\mathrm{sm}}\) be polynomial functionals of
\(\phi,\phi_b\) with compact smooth kernels. The state domain is
\begin{equation}\label{eq:vol2-native-cap-domain}
 \mathcal H^F_{\mathrm{in}}
   =\operatorname*{colim}_{N\geq\max(F\cup\{1\})}
       \left(\mathcal P_{\mathrm{sm}}\otimes_{\mathbb C}
                       Y_F[[x_1,\ldots,x_N]]\right).
\end{equation}
The series are completed by jet weight, and each weight coefficient is
polynomial in \(Y_F\). The tensor with smooth functionals is algebraic.
Completion precedes the union over the finite set of jet variables.

\begin{proposition}\label{prop:vol2-native-cap-complex}
On the domain \eqref{eq:vol2-native-cap-domain}, the normal-ordered
quantization \(Q_Y+\Omega^F_{\mathrm{in}}\) of
\eqref{eq:vol2-native-incoming-charge} squares to zero.
On the subspace of junction-jet functionals alone it is
\begin{equation}\label{eq:vol2-native-cap-differential}
                    D_G=Q_Y+\sum_{n\in F}x_nc_{n-1}.
\end{equation}
\end{proposition}

\begin{proof}
At a fixed jet stage, differentiation of a series lowers weight by a
bounded amount. Contraction with a smooth kernel takes finitely many
derivatives at the junction, while its remaining contractions restrict
smooth kernels to diagonals. The operator is therefore defined on the
stated domain.

The charge has no \(\phi_b\)-dependent coefficient. Among its terms,
only the \(a_b\)-momentum can contract with a \(\phi\)-coefficient.
Thus a commutator has at most one such contraction, and equals the
quantization of the relative Poisson bracket. The classical calculation
above gives
\(Q_YS^F_{\mathrm{in}}+\tfrac12
\{S^F_{\mathrm{in}},S^F_{\mathrm{in}}\}_{\mathrm{in}}=0\).
There is no higher-contraction correction to its quantization.
On a junction-jet functional, every momentum term has
\(\delta/\delta\phi_b\) on the right and vanishes. The endpoint
line-ghost multiplication remains, giving
\eqref{eq:vol2-native-cap-differential}.
\end{proof}

The larger completion in all jet variables together with smooth
functionals fails this domain test. Choose \(g\) smooth and compactly
supported, equal to \((1-z)^{-1}\) near zero. The sequence
\(u_M=\phi_b(g)\sum_{n=1}^M x_n\) converges in that larger jet-weight
completion. The scalar part of its quantum image is
\[
 \left.\Omega^F_{\mathrm{in}}u_M\right|_{\phi=\phi_b=0}
      =\pi\kappa\sum_{n=1}^M
                      \frac{\partial_z^ng(0)}{(n-1)!}
      =\pi\kappa\frac{M(M+1)}2.
\]
It does not stabilize even at output weight zero. This is why
\eqref{eq:vol2-native-cap-domain} retains a finite jet stage for each
mixed smooth state.

The incoming zero-momentum Lagrangian has constant flat prequantum section
away from the junction. A normalized exponential cap is \(e^f\), where
\(f\) is a polynomial in finitely many \(x_n\), has coefficients in
\(Y_F^0\), and has zero constant term in the jets. Equation
\eqref{eq:vol2-native-cap-differential} says that it is closed exactly when
\[
 f=-\sum_{n\in F}\frac{B_nx_n}{\kappa n}+h(x),
 \qquad h\in\mathbb C[B_0][x_1,x_2,\ldots].
\]
Indeed the coefficient of \(c_{n-1}\) is
\(\kappa n\partial_{B_n}f+x_n\).
Rotation gives \(x_n\) weight \(-n\), while \(B_n\) has weight \(n\).
Every nonconstant monomial of \(h\) has negative weight.
The unique rotation-invariant cap in this declared exponential class is
therefore
\begin{equation}\label{eq:vol2-native-gaussian-cap}
                G_F=\exp\left(-\sum_{n\in F}\frac{B_nx_n}{\kappa n}\right).
\end{equation}
The constant functional \(1\) is not closed when \(F\ne\varnothing\).

Gluing a linear junction functional \(\sum_ny_nx_n\) by the incoming
BV pairing gives \(\sum_ny_ne_{n-1}\) on regular source observables.
For \eqref{eq:vol2-native-gaussian-cap}, this is exactly the endpoint
correction in \eqref{eq:vol2-native-corrected-source}.
The sign is essential: the opposite exponential leaves twice the bare
endpoint defect. Nonlinear junction functionals leave point observables
after contraction and require a larger observable domain.

On the degreewise oscillator-weight product, the all-jet expression
\[
 G=\exp\left(-\sum_{n>0}\frac{x_nB_n}{\kappa n}\right)
\]
is well-defined: each weight coefficient is a finite polynomial in the
line coefficients. Define its polynomial descendant image
\(\mathsf P_\kappa=G(\mathbb C[x_1,x_2,\ldots]\otimes Y_\infty)\).
Then
\[
 D_G(Gf)=GQ_Yf,
 \qquad
 (\kappa n\partial_{x_n}+B_n)(Gf)=G\kappa n\partial_{x_n}f.
\]
Multiplication by \(G\) is a chain isomorphism onto this image, with
inverse \(G^{-1}\). Composing it with
\eqref{eq:vol2-native-regular-period} gives the regular incoming cap map,
sending \(1\) to \(G\) and retaining every line coefficient and
separated current. It identifies the cap image. It does not identify the
whole incoming state completion with a polynomial Fock module.

\section{Formal states and their Laurent fields}

The preceding completions control different variables. Completing the
coefficient ring introduces a further distinction, even for the constant
Heisenberg theory. Fix a numerical \(\kappa\ne0\), specialize the quantum
parameter to one, and introduce an independent formal parameter
\(\varepsilon\). In this section put
\(R_\varepsilon=\mathbb C[[\varepsilon]]\) and
\(R_{\varepsilon,r}=R_\varepsilon/(\varepsilon^r)\).
The parameter \(\varepsilon\) is neither a spacetime coordinate nor the
quantum parameter used to define the current commutator.

Let \(H=\mathbb C[x_1,x_2,\ldots]\) be the polynomial vacuum space,
with \(J=x_1\) and translation \(Tx_n=nx_{n+1}\).
Its current has singular product
\(Y(J,z)J=\kappa z^{-2}+\text{regular terms}\).
Every finite-precision state space
\(H\otimes R_{\varepsilon,r}\) has the usual Laurent fields.
The compatible inverse limit of the state spaces is \(H[[\varepsilon]]\).

\begin{proposition}\label{prop:vol2-native-completed-state-pole}
There is a state in \(H[[\varepsilon]]\) whose field on \(J\) has
no pole bound independent of coefficient precision. Consequently the
finite-precision fields do not take all completed state pairs into
\(H[[\varepsilon]]((z))\).
\end{proposition}

\begin{proof}
The state
\[
 a=\sum_{n\geq0}\frac{(-\varepsilon)^n}{(n+1)!}T^nJ
\]
is well-defined coefficientwise. Translation covariance gives
\(Y(T^nJ,z)=\partial_z^nY(J,z)\). Therefore its singular value is
\begin{equation}\label{eq:vol2-native-completed-pole}
             Y(a,z)J\big|_{\mathrm{sing}}
                =\kappa\sum_{n\geq0}\varepsilon^nz^{-n-2}.
\end{equation}
Modulo \(\varepsilon^r\), the pole has order \(r+1\).
All its coefficients are forced by the compatible reductions. An element
of the displayed uniform Laurent space has only finitely many negative
powers of \(z\), so it cannot have these reductions.
\end{proof}

This state is excluded by a fixed finite-weight restriction. Thus a
finite-weight state core and the total formal state completion answer
different questions. A natural field recipient for the latter is the
restricted Laurent module
\begin{equation}\label{eq:vol2-native-restricted-laurent}
 \mathcal L_\varepsilon(M)
 =\varprojlim_r(M/\varepsilon^rM)((z))
 =\left\{\sum_{j\in\mathbb Z}m_jz^j:
   \forall r\ \exists N_r\quad
   m_j\in\varepsilon^rM\text{ for }j<-N_r\right\},
\end{equation}
where \(M\) is separated and complete. The equality follows by lifting
each compatible coefficient and applying the finite pole bound at each
precision. Formula \eqref{eq:vol2-native-completed-pole} belongs to this
module. Coefficientwise Wick operations extend here because, at every
precision, they are the original finite-precision operations. This is a
formal recipient; it asserts no numerical convergence of the series.

The same distinction occurs in a separated Wilson junction. At
\(\kappa=1\), take rank-one coefficient modules with zero differential
and
\[
 B'_0=0,\quad B'_1=\varepsilon,\quad B'_n=0\ (n>1),
 \qquad B''_0=1,\quad B''_n=0\ (n>0),\qquad c'_j=c''_j=0.
\]
They satisfy the line relation \([d,B_n]=nc_{n-1}\).
The separated scalar twist is \(s(z)=\varepsilon/z\); the two
connections are
\[
                 \nabla_0=d,\qquad
                 \nabla_1=d-\varepsilon z^{-2}dz.
\]
At precision \(r\), the horizontal comparison and its inverse are
\begin{equation}\label{eq:vol2-native-formal-gauge}
 G_r(z)=\sum_{n=0}^{r-1}\frac{(-\varepsilon)^n}{n!z^n},
 \qquad G_r(z)^{-1}=\sum_{n=0}^{r-1}\frac{\varepsilon^n}{n!z^n}.
\end{equation}
Their product is one by the binomial identity, and
\(G'_r-\varepsilon z^{-2}G_r=0\).

Over \(R_\varepsilon((z))\) there is no nonzero horizontal comparison.
If \(g=\sum_{j\geq j_0}g_jz^j\ne0\), with \(g_{j_0}\ne0\), then
the coefficient of \(z^{j_0-2}\) in
\(g'-\varepsilon z^{-2}g\) is
\(-\varepsilon g_{j_0}\ne0\).
A current-linear endomorphism of the polynomial Fock factor is scalar:
its value on the vacuum is killed by every positive mode, hence is a
constant, and negative modes generate all other states. It follows that
\[
 \operatorname{Hom}^{\mathrm{flat,current}}_{R_\varepsilon((z))}=0,
 \qquad
 \varprojlim_r
 \operatorname{Hom}^{\mathrm{flat,current}}_{R_{\varepsilon,r}((z))}
       =R_\varepsilon.
\]
The comparison \(e^{-\varepsilon/z}\) does exist in
\eqref{eq:vol2-native-restricted-laurent}. These are comparisons of
separated flat current families. Extending them through coincidence
requires the supported collision operation as additional data.

There is a useful sufficient bound. If all twist coefficients lie in
an ideal whose \(m\)-th power is zero, and their largest inverse-separation
degree is \(k\), the exponential has pole order at most \(k(m-1)\).
The exponent must be uniform. In \eqref{eq:vol2-native-formal-gauge},
the nilpotence exponent grows with \(r\), and so does the pole order.

\section{Normal support after coefficient completion}

The analogous support issue is visible on one normal coordinate \(u\).
For this section abbreviate \(R=R_\varepsilon\), \(R_r=R_{\varepsilon,r}\),
and use right modules over the Weyl algebra
\[
 D_R=R\langle u,\partial_u\rangle/
                  (\partial_u u-u\partial_u-1).
\]
The parameter \(\varepsilon\) commutes with both operators.
The principal-part module and its density action are
\begin{equation}\label{eq:vol2-native-normal-principal-parts}
 P=R[u,u^{-1}]/R[u]=\bigoplus_{m\geq0}R\beta_m,
 \qquad
 \beta_0u=0,\quad \beta_mu=\beta_{m-1}\ (m>0),\quad
 \beta_m\partial_u=(m+1)\beta_{m+1}.
\end{equation}
Right differentiation is minus ordinary differentiation of the density.
In particular \(P=D_R/uD_R\), with right-module resolution
\[
             0\longrightarrow D_R\xrightarrow{L_u}D_R
                     \longrightarrow P\longrightarrow0.
\]
Here \(L_u\) is left multiplication by \(u\), and is right linear.

The coefficient completion is the restricted mode product
\begin{equation}\label{eq:vol2-native-normal-completion}
 \widehat P=
 \left\{\sum_{m\geq0}a_m\beta_m:
    a_m\in R,\quad
    \forall r\text{ all but finitely many }a_m\in\varepsilon^rR\right\}.
\end{equation}
It is derived complete because the reductions have surjective transitions
and \(P\) is scalar free. Its ordinary normal support has changed.

\begin{proposition}\label{prop:vol2-native-normal-open}
For \(Q=\widehat P/P\), both \(u\) and \(\varepsilon\) act invertibly,
and
\[
 Q\ne0,\qquad \widehat P[u^{-1}]=Q,
 \qquad Q\otimes_R^{\mathbb L}R_r=0.
\]
Thus \(P_r\) has no lift that is both derived complete with those
compatible reductions and ordinarily supported at \(u=0\).
\end{proposition}

\begin{proof}
Lowering the modes is surjective on \(\widehat P\). An element killed
by \(u^N\) has no modes above \(\beta_{N-1}\), so the union of these
kernels is exactly \(P\). Localization is therefore the quotient \(Q\),
on which \(u\) is invertible. The element
\(p=\sum_{m\geq0}\varepsilon^m\beta_m\) has nonzero image there and
satisfies \(pu=\varepsilon p\).

Subtracting a finite representative modulo \(\varepsilon\) and dividing
its remaining coefficients by \(\varepsilon\) proves surjectivity of
\(\varepsilon\) on \(Q\). If \(\varepsilon x\) has finite mode support,
so does \(x\), since \(R\) has no \(\varepsilon\)-torsion. This proves
injectivity. Derived reduction is the cone of multiplication by
\(\varepsilon^r\), and hence vanishes on \(Q\).
Finally, a derived complete lift of the compatible tower is its derived
inverse limit \(\widehat P\), whose localization is nonzero.
\end{proof}

The scalar vacuum \(R\) still exists. The proposition concerns its
normal coefficient, and the perfect test \(D_R\) detects the nonzero
open quotient by \(\operatorname{Hom}_{D_R}(D_R,Q)=Q\).

For a complex \(M\), write
\[
 M_r=M\otimes_R^{\mathbb L}R_r,
 \qquad \Lambda M=R\!\varprojlim_r M_r.
\]
Call \(M\) \emph{derived adically supported} at \(u=0\) when
\begin{equation}\label{eq:vol2-native-derived-support}
                         M_r[u^{-1}]\simeq0
                         \quad\text{for every }r.
\end{equation}
Open restriction is flat localization of the underlying
\(R[u]\)-module with its induced differential action. It does not treat
\(u\) as central in \(D_R\).
Every cohomology module of every derived reduction is tested in
\eqref{eq:vol2-native-derived-support}. For a scalar-flat module this is
equivalent to
\[
 \forall v\ \forall r\ \exists N:
                         vu^N\in\varepsilon^rM.
\]
For a general module the displayed condition controls only ordinary
quotients, and misses the negative cohomology of derived reduction.

Here is an explicit instance of that difference. Let \(M\) consist of
families in
\(\prod_{n\geq1,m\geq0}(R/\varepsilon^n)\beta_m^{(n)}\)
whose coefficients are nonzero modulo \(\varepsilon^r\) at only
finitely many pairs \((n,m)\), for each \(r\). Give each summand the
normal action \eqref{eq:vol2-native-normal-principal-parts}.
Division of coefficients shows that
\[
 M/\varepsilon^rM
       =\bigoplus_{n\geq1}P/\varepsilon^{\min(n,r)}P.
\]
These quotients are supported, and their inverse limit is \(M\).
The negative cohomology of \(M_r\) is
\(T_r=M[\varepsilon^r]\), with transition multiplication by
\(\varepsilon\). Its inverse limit is zero by separatedness. Moreover,
for \(y_r\in T_r\), the convergent series
\(x_r=\sum_{j\geq0}\varepsilon^jy_{r+j}\) belongs to \(T_r\) and
satisfies \(x_r-\varepsilon x_{r+1}=y_r\). Thus
\(\varprojlim^1T_r=0\), proving derived completeness of \(M\).

Nevertheless, in \(N=M[u^{-1}]\) the elements
\begin{equation}\label{eq:vol2-native-torsion-witness}
 w_r=\sum_{n\geq r}\varepsilon^{n-r}\beta_{n-1}^{(n)}
\end{equation}
give nonzero compatible classes in \(N[\varepsilon^r]\).
Indeed \(\varepsilon^rw_r=0\), and
\(w_r-\varepsilon w_{r+1}=\beta_{r-1}^{(r)}\) vanishes after
localization. No power of \(u\) kills \(w_r\): choose an index
\(n\geq r\) larger than that power. Hence the compatible class in
\(\varprojlim_rN[\varepsilon^r]\) is nonzero, and
\(H^{-1}(\Lambda N)\ne0\). In fact the Milnor exact sequence in
degree \(-1\) identifies these two groups, since the reductions have
no degree \(-2\) cohomology. This proves that the complete module
\(M\) fails \eqref{eq:vol2-native-derived-support} despite its
supported ordinary quotients.

\section{Reconstruction on the normal Weyl carrier}

The derived support condition has a positive reconstruction theorem.
It applies to normal coefficients; identifying them with all native
fields or their physical fusion requires additional comparisons.
Let \(\mathcal G_\varepsilon\) be the category of derived complete
scalar complexes, and let \(\mathcal C_\varepsilon\) be the category
of derived complete right \(D_R\)-complexes satisfying
\eqref{eq:vol2-native-derived-support}.

First, completion preserves every derived reduction:
\begin{equation}\label{eq:vol2-native-completion-reduction}
                         (\Lambda X)_r\simeq X_r.
\end{equation}
To see this without a boundedness assumption, tensor the inverse limit
with the perfect complex \([R\xrightarrow{\varepsilon^r}R]\).
For \(n\geq r\), the resulting two-term model of
\(X_r\otimes_R^{\mathbb L}R_n\) has zero differential between its
two copies of \(X_r\). The transitions are multiplication by
\(\varepsilon\) on the negative copy and identity on the other copy.
The first tower has zero homotopy limit: its shift on the product has
\(r\)-th power zero, so \(1-\mathrm{shift}\) has a finite geometric-series
inverse. The second tower is constant. This proves
\eqref{eq:vol2-native-completion-reduction} with its natural map.
The fibre of completion is an inverse limit of multiplication by
\(\varepsilon\), hence has invertible \(\varepsilon\). Maps from such
an object to a complete object vanish, by computing them into its
reductions. Consequently \(\Lambda\) is an exact reflection onto
complete complexes.

Localization commutes with scalar reduction. Therefore
\eqref{eq:vol2-native-derived-support} can also be tested by
\(M_1[u^{-1}]\simeq0\), or by \(\Lambda(M[u^{-1}])\simeq0\).
Indeed the first condition says that \(\varepsilon\) is invertible on
the open restriction, which is precisely the kernel of completion.

\begin{theorem}\label{thm:vol2-native-normal-reconstruction}
The functors
\[
 I(E)=\Lambda(P\otimes_R^{\mathbb L}E),\qquad
 K(M)=R\operatorname{Hom}_{D_R}(P,M)
\]
give inverse equivalences
\(\mathcal G_\varepsilon\simeq\mathcal C_\varepsilon\).
No scalar flatness or cohomological boundedness is required.
\end{theorem}

\begin{proof}
Over a commutative complex algebra \(B\), let \(V\) be a right Weyl
module with locally nilpotent \(u\). Define the finite Taylor projector
\[
 \pi(v)=\sum_{j\geq0}\frac{(-1)^j}{j!}(vu^j)\partial_u^j.
\]
The identity \(\partial_u^ju=u\partial_u^j+j\partial_u^{j-1}\)
gives \(\pi(v)u=0\) by consecutive cancellation.
The binomial identity gives the inverse formula
\[
                    v=\sum_{m\geq0}
                         \frac{\pi(vu^m)\partial_u^m}{m!}.
\]
Thus \(P_B\otimes_B\ker u\to V\), sending
\(\beta_m\otimes e\) to \(e\partial_u^m/m!\), is surjective.
Applying the largest power of \(u\) in a finite relation extracts
its highest coefficient, so the map is injective. It also proves
surjectivity of \(u\) on \(V\).

The resolution of \(P_B\) computes
\[
             K_B(V)=[V\xrightarrow{u}V]_{[0,1]}.
\]
This is a complex of scalar modules; its differential is not claimed
to be right Weyl linear. For a complex with supported cohomology,
the long exact sequence identifies
\(H^qK_B(V)=\ker(u:H^qV\to H^qV)\), since \(u\) is surjective
on that cohomology. Scalar freeness of \(P_B\) now proves that
the evaluation counit is an isomorphism on all cohomology groups.
The unit is an equivalence because lowering on \(P_B\) is split
surjective with kernel \(B\beta_0\).

The finite free resolution makes \(K\) commute with derived scalar
reduction and completion. Tensor--Hom adjunction and the reflection
property give \(I\dashv K\) on complete objects. At precision \(r\),
the unit and counit are the equivalences just proved over \(B=R_r\).
Their limits give the assertion.
\end{proof}

For example, this equivalence supplies an intrinsic supported product.
It need not be the ambient product. The latter fails already for
\(W=\prod_{i\geq0}\widehat P\). It is scalar flat and complete, and
\(W_1=\prod_iP_1\). The element
\(e=(\beta_0,\beta_1,\ldots)\) has a nonzero \(i=N\) coordinate
after multiplication by \(u^N\), for every \(N\). Thus \(W_1\) is
not supported. The intrinsic product is instead
\(I(\prod_iR)\): at each coefficient precision it imposes a mode bound
uniform over \(i\). The example concerns all indices; every finite
product is supported.

There is also an actual geometric operation compatible with the
equivalence. Let \(a:\mathbb A_x^1\times\mathbb A_y^1\to\mathbb A_t^1\)
be addition, \(t=x+y\), and define
\[
                   M*_+N=\Lambda a_+
                             (M\boxtimes_R^{\mathbb L}N),
\]
where \(a_+\) is right differential-module direct image, computed by
the relative de Rham complex. For complete scalar complexes put
\(E\widehat\otimes_R F=\Lambda(E\otimes_R^{\mathbb L}F)\).

\begin{proposition}\label{prop:vol2-native-normal-convolution}
There is a natural equivalence
\[
                    I(E)*_+I(F)\simeq
                    I(E\widehat\otimes_R F).
\]
It is compatible with associativity, symmetry and the convolution unit
\(\widehat P\). Relative convolution over an algebra \(I(B)\)
corresponds to \(I(\Lambda(E\otimes_B^{\mathbb L}F))\), with its
specified right and left module actions.
\end{proposition}

\begin{proof}
At precision \(r\), change variables to \(t=x+y\), \(s=y\).
Since \(\partial_x=\partial_t\) and
\(\partial_y=\partial_t+\partial_s\), the density modes transform by
\[
 \beta_i(x)\beta_j(y)\longmapsto
       \sum_{k=0}^j\binom{i+k}{k}
                 \beta_{i+k}(t)\beta_{j-k}(s).
\]
The relative de Rham differential on \(P_s\), in degrees \(-1,0\),
is \(\partial_s\); it raises the mode by one with coefficient \(m+1\).
It is split injective over \(R_r\), with cokernel \(R_r\beta_0\).
Projection to that cokernel sends the displayed expression to
\(\binom{i+j}{i}\beta_{i+j}(t)\). This remains a contraction after
tensoring with any derived scalar ground.
The three-input coefficients agree by the multinomial identity.
Addition and iterated relative de Rham direct image provide the
coherent comparisons, symmetry and unit.

Scalar reduction commutes with this construction. Equation
\eqref{eq:vol2-native-completion-reduction} then gives the completed
comparison. Finally its two-sided bar has the actual action maps and
terms \(E\widehat\otimes B^{\widehat\otimes n}\widehat\otimes F\).
Its realization is the completed relative tensor in the statement,
because completion is a monoidal localization: its kernel of
\(\varepsilon\)-invertible objects is a tensor ideal.
\end{proof}

This additive normal convolution retains the coordinate \(t\).
The balanced Weyl tensor consumes the normal coordinate instead:
\[
 R\operatorname{Hom}_{D_R}(P,P)=R,
 \qquad P\otimes_{D_R}^{\mathbb L}P^t=R[1].
\]
Here transpose sends \(u\) to \(u\) and \(\partial_u\) to
\(-\partial_u\). The same lowering operator computes both formulas,
in degrees \(0,1\) for Hom and \(-1,0\) for tensor.
Neither operation alone is the physical boundary fusion constructed
from separated insertions and their collision maps.

Nor is scalar flatness preserved by every relative tensor. For
\(B=\widehat{R[t]}^{\,(\varepsilon)}\), take the two rank-one modules
with \(t=0\) and \(t=\varepsilon\). The free resolution for the first
gives their relative tensor as
\([R\xrightarrow{\varepsilon}R]_{[-1,0]}\simeq R/(\varepsilon)\).
At precision \(r\), both its kernel and cokernel are nonzero.
Derived normal support retains both; ordinary scalar reduction loses
one of them.

\section{The interacting comparison}

Return to a nonconstant Hamiltonian operator \(H\).
The classical master identity constructs its local bulk action.
The Gaussian calculation constructs a quantum action on supported
boundary observables, a charged junction and an incoming cap on their
declared domains. Extending all of them to the interacting theory is
a coupled problem.

To formulate its perturbative obstruction, suppose a regulator and
a local subtraction prescription give a BV complex of bulk, boundary
and defect functionals, with the boundary terms included in its
differential. Let \(S_0\) be its classical solution and let
\(s=Q+\{S_0,-\}\). Write a proposed quantum action as
\(S=S_0+\sum_{n\geq1}\hbar^nS_n\).
In a regularized BV model its equation is
\begin{equation}\label{eq:vol2-native-joint-qme}
                 QS+\tfrac12\{S,S\}+\hbar\Delta S=0.
\end{equation}
At order \(n\), the equation for the degree-zero correction is
\[
 sS_n=-\mathcal O_n,
 \qquad
 \mathcal O_n=\Delta S_{n-1}
       +\tfrac12\sum_{i=1}^{n-1}\{S_i,S_{n-i}\}.
\]
The BV identities and the lower-order equations give
\(s\mathcal O_n=0\), by taking the order-\(n\) term of the identity
that the differential twisted by \(S\) annihilates its curvature.
Thus the obstruction is a class in degree one of this joint local
complex. A class in the bulk-only complex does not test endpoint
counterterms or the boundary Ward identities.

The required local prescription must also respect restrictions to
nested regions. Its annulus amplitudes must sew with the chosen cap,
and collision extensions must obey the same current and defect Ward
identities. The forest contractions of
Chapter~\ref{ch:vol2-supported-compositions} provide explicit Gaussian
homotopies. They do not solve the interacting master equation on all
these supports.

An algebraic Hochschild action poses a different comparison. To
identify it with the physical action, one needs a map from the
renormalized bulk observables to the coherent operations on the
boundary category, preserving the unit, translation, collision maps
and their homotopies. If a centre represents those operations, the map
must factor through that representing object. Equivalence further
requires that the chosen boundary objects detect and reconstruct the
bulk. Neither a shared current bracket nor a solution of the master
equation supplies these maps by itself.

The Gaussian normal contraction shows why the operation is the relevant
datum: a field removed by linear projection can still act through its
contraction with a boundary insertion. Formal support and Laurent
completion then specify where that action lives. The general interacting
bulk-to-boundary evaluator, and its identification with a universal
holomorphic centre, remain open beyond the constructions proved here.
